%% ---------------------------------------------------------------------------------------
\section{Double descent}
\label{sec:dd}

\paragraph{Protocol.} Double descent is easiest to see with few samples and label noise
\cite{belkin,nakkiran}, so we use $n=4\,000$ CIFAR-10 training images (ImageNet trunk, PCA-64), with
clean labels and with 20\% of the labels replaced by a random wrong class, and the full 10\,000-image
test set. For every width $H\in\{1,2,4,\dots\}$ we train one MLP $64\!\to\!H\!\to\!10$ with
full-batch Adam for 3\,000 steps and no regularisation, and then grow a corrector tree that uses
\emph{that same MLP as its root} and binary nodes of width $H$, until it interpolates. The two curves in
Fig.~\ref{fig:dd} therefore differ only by the children.

\begin{figure}[t]
\centering
\includegraphics[width=0.8\linewidth]{double_descent.pdf}
\caption{Model-wise double descent on CIFAR-10 ($n=4\,000$). Blue: the root MLP alone; orange: the
tree built on that root, which always reaches 0\% training error. Dashed line: the smallest width at
which the MLP itself interpolates.}
\label{fig:dd}
\end{figure}

\paragraph{The root network shows a textbook double descent.} With 20\% noise (Fig.~\ref{fig:dd},
left) the MLP's test error falls from 75.7\% ($H=1$) to 33.2\% ($H=8$), the classical sweet spot, then
\emph{rises} to 53.8\% at $H=64$ -- precisely the first width whose training error is zero, where it has
memorised all 815 corrupted labels -- and then falls again monotonically to 39.0\% at $H=4\,096$
(307k parameters, $7.7\times$ more than $nK=40$k). The test cross-entropy peaks at the same width
(7.7 vs 1.06 at $H=8$). With clean labels (right) the curve has the same shape but a smaller peak at a
smaller width: 31.4\% at $H=8$, 37.7\% at the first interpolating width $H=32$, and 30.2\% at
$H=2048$ -- here the over-parameterised branch ends \emph{below} the classical optimum. Noise moves the
interpolation threshold to the right (noisy labels need more capacity to memorise) and makes the
peak higher (memorising them at the threshold forces the least smooth interpolant).

\paragraph{The tree never sees the classical regime.} For $H$ at or above the threshold the root
interpolates on its own, no children are grown, and the two curves coincide -- including the peak,
so the tree's own error-versus-width curve has the double-descent shape too, with its peak at the
width at which the \emph{root} can just interpolate. Below the threshold the children force
interpolation and the tree's error is almost flat (noisy: 47.2--48.5\% for $H=2$--32; clean:
34.7--39.7\% for $H=2$--16), close to the error of the MLP just past its peak. Growing children is a
second way to cross the interpolation threshold: instead of widening one network we add small ones
until the data are memorised. Fig.~\ref{fig:ddp} places the trees on the parameter axis: for $H\ge2$ they hold 12k--95k
parameters, one to two orders of magnitude more than their roots and as many as MLPs of width $H=256$--1024, and their
test error is within a few points of those MLPs on the over-parameterised branch (noisy: 47--48\% vs
40--44\%) -- not that of the small root they started from (33\% at $H=8$).

\begin{figure}[t]
\centering
\includegraphics[width=0.8\linewidth]{double_descent_params.pdf}
\caption{The same runs against the number of parameters (log scale). Trees (squares, root width
$H$ below the threshold) sit with the over-parameterised MLPs, not with their own small roots.}
\label{fig:ddp}
\end{figure}

\paragraph{Children help an under-fitted root and hurt a well-fitted one.} At $H=1$--2 the root is
far from the data and the children act as a boosting stage: with clean labels the tree on a
width-1 root has 48.1\% error instead of 74.4\% (2\,916 test predictions fixed, 290 broken). From
$H=4$ on, the root is close to its best test error and the children mostly memorise its residual --
with noise, mostly the corrupted labels -- and break more than they fix (e.g.\ $H=8$, noisy: 296
fixed, 1\,788 broken). Every tree memorises 100\% of the corrupted labels; the roots below the threshold memorise
4--11\% of them ($H\le16$), rising to 39\% at $H=32$.

\paragraph{Where does the label noise go?} We grew a tree on 10\,000 CIFAR-10 images with 20\%
corrupted labels (ResNet trunk, $h=8$; 25 nodes, 100\% training accuracy, test accuracy 56.1\%$\to$50.7\%).
Fig.~\ref{fig:noise} follows the corrupted labels through it. The root memorises 47\% of them by
itself, but its training errors are 55\% corrupted labels against 20\% in the data -- the root learns
the clean structure and leaves most of the noise to its children, whose first level is therefore
mostly a \emph{noise memoriser}. Deeper nodes, which correct the detectors' own false alarms and
misses, receive clean labels again (22\% at depth 2, none at depth 3). In this sense a tree separates
signal (root) from memorisation (children) by construction, and \S\ref{sec:min} shows that removing
children is the simplest way to recover the root's generalisation.

\begin{figure}[h]
\centering
\includegraphics[width=0.5\linewidth]{noise_depth.pdf}
\caption{Fraction of corrupted labels among the samples that each part of the tree must fire on
(CIFAR-10, 10k images, 20\% symmetric noise).}
\label{fig:noise}
\end{figure}
